\begin{name}
	{\tenchude}{\tendethi}{LỚP TOÁN THẦY PHÁT}{\thoigian}
\end{name}
\setcounter{ex}{0}\setcounter{bt}{0}
\Opensolutionfile{ans}[ans/ans-2-TT-19-ChuyenHaLong-QuangNinh-L2-23]
\begin{ex}%[Thi thử TN - Chuyên Hạ Long-Quảng Ninh-L2-23]%[Duong Xuan Loi-EX6]%[1D2Y2-1]
	Tập hợp $A$ gồm $2022$ phần tử. Hỏi có bao nhiêu tập con của $A$ gồm đúng $4$ phần tử?
	\choice
	{$\mathrm{A}_{2022}^4$}
	{\True $\mathrm{C}_{2022}^4$}
	{$2022\cdot4!$}
	{$\dfrac{2023}{4!}$}
	\loigiai{
		Số tập con của $A$ gồm đúng $4$ phần tử là $\mathrm{C}_{2022}^4$.}
\end{ex}

\begin{ex}%[Thi thử TN - Chuyên Hạ Long-Quảng Ninh-L2-23]%[Duong Xuan Loi-EX6]%[1D3Y3-3]
	Cho cấp số cộng $(u_n)$ có $u_3=5;u_5=9$. Tính $u_{15}$. 
	\choice
	{$u_{15}=-29$}
	{$u_{15}=4$}
	{$u_{15}=45$}
	{\True $u_{15}=29$}
	\loigiai{
		Ta có $\heva{&u_3=5\\&u_5=9}\Leftrightarrow\heva{&u_1+2d=5\\&u_1+4d=9}\Leftrightarrow\heva{&u_1=1\\&d=2.}$ \\
		Khi đó $u_{15}=u_1+14d=1+14\cdot 2=29$.}
\end{ex}

\begin{ex}%[Thi thử TN - Chuyên Hạ Long-Quảng Ninh-L2-23]%[Duong Xuan Loi-EX6]%[2D1Y1-1]
	Hàm số nào sau đây nghịch biến trên $\mathbb{R}$?
	\choice
	{\True $y=\left(\dfrac{\pi}{6}\right)^x$}
	{$y=\log_{\tfrac{1}{2}}x$}
	{$y=x^4-3x^2+2$}
	{$y=\dfrac{x+1}{x-2}$}
	\loigiai{
		Hàm số $y=\left(\dfrac{\pi}{6}\right)^x$ có cơ số $0<\dfrac{\pi}{6}<1$ nên hàm số nghịch biến trên $\mathbb{R}$.}
\end{ex}

\begin{ex}%[Thi thử TN - Chuyên Hạ Long-Quảng Ninh-L2-23]%[Duong Xuan Loi-EX6]%[2D1Y1-2]
	Cho hàm số $y=f(x)$ có bảng biến thiên như sau
	\begin{center}
		\begin{tikzpicture}
			\tkzTabInit[nocadre=false,lgt=1.2,espcl=2.5,deltacl=0.6]
			{$x$ /0.6, $y'$ /0.6, $y$ /2.5}
			{$-\infty$,$1$,$2$,$+\infty$}
			\tkzTabLine{,+,$0$,-,$0$,+,}
			\tkzTabVar{-/$-\infty$,+/$4$,-/$2$,+/$+\infty$}
		\end{tikzpicture}
	\end{center}
	Khẳng định nào dưới đây \textbf{sai}?
	\choice
	{Hàm số đồng biến trên khoảng $(5;2023)$}
	{Hàm số nghịch biến trên khoảng $(1;2)$}
	{\True Hàm số đồng biến trên khoảng $(-\infty;4)$}
	{Hàm số đồng biến trên khoảng $(2;+\infty)$}
	\loigiai{
		Khẳng đinh sai \lq\lq Hàm số đồng biến trên khoảng $(-\infty;4)$\rq\rq.}
\end{ex}

\begin{ex}%[Thi thử TN - Chuyên Hạ Long-Quảng Ninh-L2-23]%[Duong Xuan Loi-EX6]%[2D1Y4-1]
	Đồ thị hàm số nào sau đây có tiệm cận đứng.
	\choice
	{$y=\dfrac{\sqrt{1-x^2}}{x^2+1}$}
	{\True $y=\dfrac{x^2-3x+2}{x^4-1}$}
	{$y=2^x$}
	{$y=x^3+3x^2+x+1$}
	\loigiai{
		Xét hàm số $y=\dfrac{x^2-3x+2}{x^4-1}$.\\
		Tập xác định $\mathscr{D}=\mathbb{R}\setminus\{\pm 1\}$.\\
		Ta có $\lim\limits_{x\to 1^+} y=\lim\limits_{x\to 1^+}\dfrac{(x-1)(x-2)}{(x-1)(x+1)\left(x^2+1\right)}=\lim\limits_{x\to 1^+}\dfrac{x-2}{(x+1)\left(x^2+1\right)}=-\dfrac{1}{4}$.\\
		$\lim\limits_{x\to 1^-} y=\lim\limits_{x\to 1^-}\dfrac{(x-1)(x-2)}{(x-1)(x+1)\left(x^2+1\right)}=\lim\limits_{x\to 1^-}\dfrac{x-2}{(x+1)\left(x^2+1\right)}=-\dfrac{1}{4}$.\\
		Suy ra $x=1$ không là đường tiệm cận đứng.\\
		Ta có $\lim\limits_{x\to(-1)^+} y=\lim\limits_{x\to(-1)^+}\dfrac{x^2-3x+2}{(x-1)(x+1)\left(x^2+1\right)}=-\infty;\lim\limits_{x\to (-1)^-} y=\lim\limits_{x\to(-1)^-}\dfrac{x^2-3x+2}{(x-1)(x+1)\left(x^2+1\right)}=+\infty$.\\
		Suy ra $x=-1$ là đường tiệm cận đứng.}
\end{ex}

\begin{ex}%[Thi thử TN - Chuyên Hạ Long-Quảng Ninh-L2-23]%[Duong Xuan Loi-EX6]%[2D1Y4-1]
	Đồ thị hàm số $y=\dfrac{\sqrt{4-x^2}}{-x^3+7x-6}$ có bao nhiêu đường tiệm cận?
	\choice
	{$4$}
	{\True $2$}
	{$3$}
	{$1$}
	\loigiai{
		Điều kiện xác định $-2\leq x\leq 2$.\\
		$y=\dfrac{\sqrt{4-x^2}}{-x^3+7x-6}=\dfrac{\sqrt{(2-x)(2+x)}}{-(x-1)(x-2)(x+3)}=\dfrac{\sqrt{2+x}}{(x-1)(x+3)\sqrt{2-x}}$.\\
		Vậy hàm số có $2$ tiệm cận đứng $x=1; x=2$ và không có tiệm cận ngang.}
\end{ex}

\begin{ex}%[Thi thử TN - Chuyên Hạ Long-Quảng Ninh-L2-23]%[Duong Xuan Loi-EX6]%[2D1Y5-1]
	Trong các hàm số sau, hàm số nào có bảng biến thiên như sau
	\begin{center}
		\begin{tikzpicture}
			\tikzset{double style/.append style={double distance=1.5pt}}
			\tkzTabInit[nocadre=false,lgt=1.2,espcl=3.5,deltacl=0.6]
			{$x$ /0.6,$y'$ /0.6,$y$ /2}
			{$-\infty$,$2$,$+\infty$}
			\tkzTabLine{,+,d,+,}
			\tkzTabVar{-/$1$,+D-/$+\infty$/$-\infty$,+/$1$}
		\end{tikzpicture}
	\end{center}
	\choice
	{$y=\dfrac{x+1}{x-2}$}
	{$y=\dfrac{x+3}{x-2}$}
	{$y=\dfrac{x-1}{x-2}$}
	{\True $y=\dfrac{x-3}{x-2}$}
	\loigiai{
		Dựa vào bảng biến thiên ta thấy hàm số đồng biến trên $(-\infty;2)$ và $(2;+\infty)$. Chỉ có phương án $y=\dfrac{x-3}{x-2}\Rightarrow y'=\dfrac{1}{(x-2)^2}>0,\forall\neq 2$.}
\end{ex}

\begin{ex}%[Thi thử TN - Chuyên Hạ Long-Quảng Ninh-L2-23]%[Duong Xuan Loi-EX6]%[2D1Y5-3]
	Cho hàm số $f(x)=ax^3+bx^2+cx+d\,(a\neq 0)$ có bảng biến thiên như sau 
	\begin{center}
		\begin{tikzpicture}
			\tkzTabInit[nocadre=false,lgt=1.2,espcl=2.5,deltacl=0.6]
			{$x$ /0.6, $f'(x)$ /0.6, $f(x)$ /2.5}
			{$-\infty$,$1$,$2$,$+\infty$}
			\tkzTabLine{,+,$0$,-,$0$,+,}
			\tkzTabVar{-/$-\infty$,+/$3$,-/$-1$,+/$+\infty$}
		\end{tikzpicture}
	\end{center}
	Có bao nhiêu giá trị nguyên của $m$ để phương trình $f(x)=m$ có đúng $3$ nghiệm phân biệt. 
	\choice
	{\True $3$}
	{$1$}
	{$5$}
	{$4$}
	\loigiai{
		Xét tương giao giữa hai đồ thị hàm số ta có 
		\begin{center}
			\begin{tikzpicture}
				\tkzTabInit[nocadre=false,lgt=1.2,espcl=2.5,deltacl=0.6]
				{$x$ /0.6, $f'(x)$ /0.6, $f(x)$ /2.5}
				{$-\infty$,$1$,$2$,$+\infty$}
				\tkzTabLine{,+,$0$,-,$0$,+,}
				\tkzTabVar{-/$-\infty$,+/$3$,-/$-1$,+/$+\infty$}
				\draw (1.7,-2.5)--(9,-2.5);
			\end{tikzpicture}
		\end{center}
		Phương trình $f(x)=m$ có đúng $3$ nghiệm phân biệt $\Leftrightarrow-1<m<3\Rightarrow m\in\{0;1;2\}$.}
\end{ex}

\begin{ex}%[Thi thử TN - Chuyên Hạ Long-Quảng Ninh-L2-23]%[Duong Xuan Loi-EX6]%[2D1Y5-4]
	Đồ thị hàm số $y=x^3+7x+8$ cắt trục hoành tại bao nhiêu điểm có hoành độ âm. 
	\choice
	{\True $1$}
	{$2$}
	{$0$}
	{$3$}
	\loigiai{
		Xét phương trình hoành độ giao điểm của đồ thị hàm số $y=x^3+7x+8$ cắt trục hoành.\\
		$x^3+7x+8=0\Leftrightarrow x=-1$.\\
		Đồ thị hàm số $y=x^3+7x+8$ cắt trục hoành tại 1 điểm có hoành độ âm.}
\end{ex}

\begin{ex}%[Thi thử TN - Chuyên Hạ Long-Quảng Ninh-L2-23]%[Duong Xuan Loi-EX6]%[2D2Y3-1]
	Cho $\log_35=a$. Tính $\log_{243}1125$ theo $a$. 
	\choice
	{$\dfrac{3+2a}{5}$}
	{$\dfrac{3a}{5}$}
	{$\dfrac{2+3a}{4}$}
	{\True $\dfrac{2+3a}{5}$}
	\loigiai{
		Ta có $\log_{243}1125=\log_{3^5}\left(3^2{\cdot 5}^3\right)=\dfrac{1}{5}\left(\log_33^2+\log_{_3}5^3\right)=\dfrac{2+3a}{5}$.}
\end{ex}

\begin{ex}%[Thi thử TN - Chuyên Hạ Long-Quảng Ninh-L2-23]%[Duong Xuan Loi-EX6]%[2D2Y4-2]
	Tính đạo hàm của hàm số $f(x)=3^{x^3+3x}$. 
	\choice
	{$f'(x)=\left(x^2+1\right)3^{x^3+3x+1}$}
	{$f'(x)=\left(x^3+3x\right)3^{x^3+3x-1}$}
	{\True $f'(x)=\left(x^2+1\right)3^{x^3+3x+1}\ln 3$}
	{$f'(x)=\left(x^2+1\right)3^{x^3+3x}\ln 3$}
	\loigiai{
		Ta có: $f'(x)=\left(3x^2+3\right)3^{x^3+3x}\ln 3=\left(x^2+1\right)3^{x^3+3x+1}\ln 3$.}
\end{ex}

\begin{ex}%[Thi thử TN - Chuyên Hạ Long-Quảng Ninh-L2-23]%[Duong Xuan Loi-EX6]%[2D2Y6-1]
	Tập nghiệm của bất phương trình $\log_2(x-3)>0$.
	\choice
	{$(3;+\infty)$}
	{$[4;+\infty)$}
	{$[3;+\infty)$}
	{\True $(4;+\infty)$}
	\loigiai{
		Ta có $\log_2(x-3)>0\Leftrightarrow\heva{&x-3>0\\&x-3>2^0}\Leftrightarrow\heva{&x>3\\&x>4}\Leftrightarrow x>4$.\\
		Tập nghiệm của bất phương trình là $S=(4;+\infty)$.}
\end{ex}

\begin{ex}%[Thi thử TN - Chuyên Hạ Long-Quảng Ninh-L2-23]%[Duong Xuan Loi-EX6]%[2D3Y1-1]
	Cho hàm số $f(x)$ liên tục trên $\mathbb{R}$ thỏa mãn $\displaystyle\int f(x)\mathrm{\,d}x=2x^2+x+1+C,\forall x\in\mathbb{R},C$ là hằng số. Tính $f(2023)$. 
	\choice
	{$4047$}
	{$4046$}
	{\True $8093$}
	{$8092$}
	\loigiai{
		Ta có $f(x)=4x+1$. Suy ra $f(2023)=8093$.}
\end{ex}

\begin{ex}%[Thi thử TN - Chuyên Hạ Long-Quảng Ninh-L2-23]%[Duong Xuan Loi-EX6]%[2D3Y2-1]
	Trong các mệnh đề sau có bao nhiêu mệnh đề đúng?\\
	i, $\displaystyle\int\limits_a^b f(x)\mathrm{\,d}x=-\displaystyle\int\limits_a^b-f(x)\mathrm{\,d}x$ với $f(x)$ liên tục trên $[a;b]$.\\
	ii, $\displaystyle\int\limits_a^b kf(x)\mathrm{\,d}x=k\displaystyle\int\limits_a^b f(x)\mathrm{\,d}x$ với $f(x)$ liên tục trên $[a;b]$; $k\in\mathbb{R}$.\\
	iii, $\displaystyle\int\limits_a^b f(kx)\mathrm{\,d}x=k\displaystyle\int\limits_a^b f(x)\mathrm{\,d}x$ với $f(x)$ liên tục trên $[a;b]$; $k\in\mathbb{R}$.\\
	iv, $\displaystyle\int\limits_a^b f(x)\mathrm{\,d}x+\displaystyle\int\limits_b^a g(x)\mathrm{\,d}x=\displaystyle\int\limits_a^b[f(x)-g(x)]\mathrm{\,d}x$. 
	\choice
	{$1$}
	{$4$}
	{\True $3$}
	{$2$}
	\loigiai{
		Các khẳng định đúng là i, ii, iv. Khẳng định iii sai.}
\end{ex}

\begin{ex}%[Thi thử TN - Chuyên Hạ Long-Quảng Ninh-L2-23]%[Duong Xuan Loi-EX6]%[2D4Y2-1]
	Cho $z_1=3-2i$, $z_2=4+5i$. Tính $z_1+z_2$. 
	\choice
	{$z_1+z_2=7-7i$}
	{$z_1+z_2=1+7i$}
	{\True $z_1+z_2=7+3i$}
	{$z_1+z_2=12-10i$}
	\loigiai{
		Ta có $z_1+z_2=(3-2i)+(4+5i)=(3+4)+(-2+5)i=7+3i$.}
\end{ex}

\begin{ex}%[Thi thử TN - Chuyên Hạ Long-Quảng Ninh-L2-23]%[Duong Xuan Loi-EX6]%[2D4Y2-2]
	Cho số phức $z=4+5i$. Số phức đối của $z$ là 
	\choice
	{$-z=-4+5i$}
	{$-z=4-5i$}
	{$-z=-5-4i$}
	{\True $-z=-4-5i$}
	\loigiai{
		Số phức đối của $z$ là $-z=-4-5i$.}
\end{ex}

\begin{ex}%[Thi thử TN - Chuyên Hạ Long-Quảng Ninh-L2-23]%[Duong Xuan Loi-EX6]%[2D4Y2-2]
	Xác định mô-đun của số phức $z=(1-i)^6+(1+i)^3$. 
	\choice
	{$|z|=(\sqrt{2})^6+(\sqrt{2})^3$}
	{$|z|=\sqrt{82}$}
	{$|z|=5\sqrt{2}$}
	{\True $|z|=2\sqrt{26}$}
	\loigiai{
		Ta có $z=(1-i)^6+(1+i)^3=-2+10i\Rightarrow|z|=\sqrt{(-2)^2+10^2}=2\sqrt{26}$.}
\end{ex}

\begin{ex}%[Thi thử TN - Chuyên Hạ Long-Quảng Ninh-L2-23]%[Duong Xuan Loi-EX6]%[2H1Y3-2]
	Cho lăng trụ đứng $ABC.A'B'C'$ có đáy $ABC$ là tam giác vuông tại $A$, $AC=6a$, $BC=10a$, $AA'=8a$. Tính thể tích khối chóp $A'.ABC$. 
	\choice
	{$160a^3$}
	{$192a^3$}
	{$128a^3$}
	{\True $64a^3$}
	\loigiai{
		\immini{
			Vì tam giác $ABC$ vuông tại $A$ nên ta có\\ $AB=\sqrt{BC^2-AC^2}=\sqrt{(10a)^2-(6a)^2}=8a$.\\
			Diện tích tam giác $ABC$ là $S=\dfrac{1}{2}AB\cdot AC=\dfrac{1}{2}8a\cdot 6a=24a^2$.\\
			Thể tích khối chóp $A'.ABC$ là $V=\dfrac{1}{3}\cdot S\cdot AA'=\dfrac{1}{3}\cdot 24a^2\cdot 8a=64a^3$.
		}{
			\begin{tikzpicture}[scale=0.7, font=\footnotesize, line join=round, line cap=round,>=stealth]
				\path
				(0,0)coordinate (A)
				(2.5,-1.5)coordinate (B)
				(4,0)coordinate (C)
				($(A)+(0,3.5)$) coordinate (A')
				($(A')+(B)-(A)$)coordinate (B')
				($(A')+(C)-(A)$)coordinate (C')			
				;
				\draw (A)--(B)--(C)--(C')--(B')--(A')--cycle (A')--(C') (A')--(B)--(B');
				\draw[dashed] (A)--(C)--(A');	
				\foreach \p/\q in {A/180,B/-90,C/0,A'/180,B'/75,C'/0}
				\fill[black] (\p) circle (1.0pt)node[shift={(\q:3mm)}]{$\p$};
			\end{tikzpicture}
	}}
\end{ex}

\begin{ex}%[Thi thử TN - Chuyên Hạ Long-Quảng Ninh-L2-23]%[Duong Xuan Loi-EX6]%[2H1Y3-2]
	Khối lăng trụ đứng có cạnh bên bằng $4a$, mặt đáy là tứ giác có diện tích là $36a^2$. Thể tích khối lăng trụ đã cho là
	\choice
	{$48$}
	{$48a^3$}
	{\True $144a^3$}
	{$144$}
	\loigiai{
		Ta có thể tích khối lăng trụ là $V=36a^2\cdot 4a=144a^3$.}
\end{ex}

\begin{ex}%[Thi thử TN - Chuyên Hạ Long-Quảng Ninh-L2-23]%[Duong Xuan Loi-EX6]%[2H2Y2-1]
	Khối cầu có đường kính $8a$. Thể tích khối cầu đó là
	\choice
	{\True $\dfrac{256\pi a^3}{3}$}
	{$\dfrac{2048\pi}{3}$}
	{$\dfrac{2048\pi a^2}{3}$}
	{$\dfrac{256\pi}{3}$}
	\loigiai{
		Khối cầu có bán kính $r=4a$ nên có thể tích $V=\dfrac{4}{3}\pi r^3=\dfrac{256}{3}\pi a^3$.}
\end{ex}

\begin{ex}%[Thi thử TN - Chuyên Hạ Long-Quảng Ninh-L2-23]%[Duong Xuan Loi-EX6]%[2H3Y2-2]
	Trong không gian cho hệ trục $Oxyz$, điểm $A(1;0;0)$, $B(0;-2;0)$ và $C(0;0;3)$. Phương trình mặt phẳng $(ABC)$ là
	\choice
	{\True $x-\dfrac{1}{2}y+\dfrac{1}{3}z=1$}
	{$x-2y+3z=1$}
	{$\dfrac{x}{1}+\dfrac{y}{2}+\dfrac{z}{3}=1$}
	{$x-2y+3z=0$}
	\loigiai{
		Phương trình mặt phẳng $(ABC)$ là $x-\dfrac{1}{2}y+\dfrac{1}{3}z=1$.}
\end{ex}

\begin{ex}%[Thi thử TN - Chuyên Hạ Long-Quảng Ninh-L2-23]%[Duong Xuan Loi-EX6]%[2H3Y3-1]
	Trong không gian cho hệ trục $Oxyz$, đường thẳng $\Delta$ có phương trình $\dfrac{x-1}{2}=\dfrac{y+1}{3}=\dfrac{-z-1}{4}$. Đường thẳng $\Delta$ có véc-tơ chỉ phương là 
	\choice
	{$\overrightarrow{u}(1;-1; 1)$}
	{\True $\overrightarrow{u}(2; 3;-4)$}
	{$\overrightarrow{u}(2; 3; 4)$}
	{$\overrightarrow{u}(-1; 1;-1)$}
	\loigiai{
		Đường thẳng $\Delta$ có phương trình $\dfrac{x-1}{2}=\dfrac{y+1}{3}=\dfrac{-z-1}{4}\Leftrightarrow\dfrac{x-1}{2}=\dfrac{y+1}{3}=\dfrac{z+1}{-4}$.\\
		Suy ra đường thẳng $\Delta$ có véc-tơ chỉ phương là $\overrightarrow{u}(2; 3;-4)$.}
\end{ex}

\begin{ex}%[Thi thử TN - Chuyên Hạ Long-Quảng Ninh-L2-23]%[Duong Xuan Loi-EX6]%[1D2B5-2]
	Một hộp có $5$ quả cầu vàng, $7$ quả cầu đỏ, $6$ quả cầu xanh. Chọn ngẫu nhiên từ hộp $4$ quả cầu. Tính xác suất để $4$ quả cầu lấy được có đủ $3$ màu khác nhau. 
	\choice
	{$\dfrac{165}{408}$}
	{$\dfrac{35}{612}$}
	{\True $\dfrac{35}{68}$}
	{$\dfrac{225}{3060}$}
	\loigiai{
		Ta có $n(\Omega)=\mathrm{C}_{18}^4=3060$.\\
		Gọi $A$ là biến cố “Chọn ngẫu nhiên từ hộp $4$ quả cầu để $4$ quả cầu lấy được có đủ $3$ màu khác nhau”.\\
		Để chọn ngẫu nhiên từ hộp $4$ quả cầu để $4$ quả cầu lấy được có đủ $3$ màu khác nhau, ta có các trường hợp sau:\\
		$2$ quả cầu vàng, $1$ quả cầu đỏ, $1$ quả cầu xanh;\\
		$1$ quả cầu vàng, $2$ quả cầu đỏ, $1$ quả cầu xanh;\\
		$1$ quả cầu vàng, $1$ quả cầu đỏ, $2$ quả cầu xanh.\\
		Vậy $n(A)=\mathrm{C}_5^2\mathrm{C}_7^1\mathrm{C}_6^1+\mathrm{C}_5^1\mathrm{C}_7^2\mathrm{C}_6^1+\mathrm{C}_5^1\mathrm{C}_7^1\mathrm{C}_6^2=1575\Rightarrow\mathrm{P}(A)=\dfrac{n(A)}{n(\Omega)}=\dfrac{1575}{3060}=\dfrac{35}{68}$.}
\end{ex}

\begin{ex}%[Thi thử TN - Chuyên Hạ Long-Quảng Ninh-L2-23]%[Duong Xuan Loi-EX6]%[1H3B5-3]
	Cho hình lập phương $ABCD.A'B'C'D'$ có cạnh bằng $a$. Tính khoảng cách giữa hai mặt phẳng $(AB'D')$ và $(BDC')$. 
	\choice
	{$\dfrac{a}{3}$}
	{$\dfrac{a\sqrt{2}}{3}$}
	{$\dfrac{a\sqrt{3}}{2}$}
	{\True $\dfrac{a\sqrt{3}}{3}$}
	\loigiai{
		\immini{
			Ta có $(AB'D')\parallel(BDC')$ nên $$\mathrm{d}\left((AB'D'),(BDC')\right)=\mathrm{d}\left(C',(AB'D')\right)=\mathrm{d}\left(A',(AB'D')\right)=h.$$
		Do khối tứ diện $A'AB'D'$ vuông tại $A'$, dó đó 
		\[\dfrac{1}{h^2}=\dfrac{1}{(AA')^2}+\dfrac{1}{(A'B')^2}+\dfrac{1}{(A'D')^2}=\dfrac{3}{a^2}\Rightarrow h=\dfrac{a\sqrt{3}}{3}.\]
		}
		{
			\begin{tikzpicture}[scale=0.7,font=\footnotesize,line join=round,line cap=round,>=stealth]
				\path
				(0,0) coordinate (B)
				(1,1.4) coordinate (A)
				(4,0) coordinate (C)
				($(C)-(B)+(A)$) coordinate (D)	
				($(A)+(90:3.5)$) coordinate (A')
				($(B)-(A)+(A')$) coordinate (B')
				($(C)-(A)+(A')$) coordinate (C')
				($(D)-(A)+(A')$) coordinate (D')
				;
				\draw (B')--(B)--(C)--(D)--(D')--(A')--(B')--(C')--(D')--cycle (C)--(C')--(B) (C')--(D);
				\draw[dashed] (D')--(A)--(D) (A')--(A)--(B)--(D) (A)--(B');
				\foreach \p/\q in {A/160,B/-135,C/-45,D/0,A'/90,B'/180,C'/-20,D'/0}
				\fill[black] (\p)node[shift={(\q:3mm)}]{$\p$} circle (1.0pt);	
			\end{tikzpicture}
	}}
\end{ex}

\begin{ex}%[Thi thử TN - Chuyên Hạ Long-Quảng Ninh-L2-23]%[Duong Xuan Loi-EX6]%[2D1B2-1]
	Cho hàm số $f(x)$ có đạo hàm $f'(x)=(x+1)^2(x-1)(x+2)^3$. Đồ thị hàm số $f(x)$ có bao nhiêu điểm cực trị?
	\choice
	{$0$}
	{$3$}
	{$1$}
	{\True $2$}
	\loigiai{
		$f'(x)=0\Leftrightarrow\hoac{&x=-1\\&x=1\\&x=-2.}$ \\
		$f'(x)=0$ có hai nghiệm đơn là $x=1; x=-2\Rightarrow f'(x)$ đổi dấu hai lần qua hai nghiệm nên đồ thị hàm số có hai điểm cực trị.
	}
\end{ex}

\begin{ex}%[Thi thử TN - Chuyên Hạ Long-Quảng Ninh-L2-23]%[Duong Xuan Loi-EX6]%[2D1B2-3]
	Có bao nhiêu giá trị của $m$ để hàm số $y=x^3-(2m+3)x^2+\left(m^2+2m\right)x+2$ đạt cực đại tại $x=0$?
	\choice
	{$0$}
	{Vô số}
	{\True $1$}
	{$2$}
	\loigiai{
		$y'=3x^2-2(2m+3)x+m^2+2m$.\\
		Hàm số đạt cực đại tại $x=0\Rightarrow y'(0)=0\Leftrightarrow m^2+2m=0\Leftrightarrow\hoac{&m=0\\&m=-2.}$ \\
		Với $m=0$.\\
		Ta có $y'=3x^2-6x$. Lập bảng biến thiên suy ra $x=0$ là điểm cực đại.\\
		Với $m=-2$.\\
		Ta có $y'=3x^2+2x$. Lập bảng biến thiên suy ra $x=0$ là điểm cực tiểu.\\
		Vậy $m=0$ thỏa mãn yêu cầu bài toán.}
\end{ex}

\begin{ex}%[Thi thử TN - Chuyên Hạ Long-Quảng Ninh-L2-23]%[Duong Xuan Loi-EX6]%[2D1B3-1]
	Biết hàm số $y=\sqrt{\sin^2x+2\sin x+2}+\sqrt{\sin^2x-4\sin x+13}$ đạt giá trị nhỏ nhất khi $\sin x=\dfrac{a}{b} (\dfrac{a}{b}$ là phân số tối giản; $a$, $b$ là các số nguyên; $b>0)$. Tính $a^3+b^2$. 
	\choice
	{\True $15$}
	{$17$}
	{$65$}
	{$3$}
	\loigiai{
		Đặt $\sin x=t\Rightarrow t\in[-1;1]$.\\
		$y=\sqrt{t^2+2t+2}+\sqrt{t^2-4t+13}=\sqrt{(t+1)^2+1}+\sqrt{(t-2)^2+9}=MA+MB.$\\
		Với $M(t;0),A(-1;-1),B(2;3)$.
		\begin{center}
			\begin{tikzpicture}[scale=1, font=\footnotesize, line join=round, line cap=round,>=stealth]
				\def\xmin{-2} \def\xmax{3}
				\def\ymin{-1.5} \def\ymax{4}				
				\draw[->] (\xmin,0)--(\xmax,0) node [below]{$x$};
				\draw[->] (0,\ymin)--(0,\ymax) node [left]{$y$};
				\node at (0,0) [below right]{$O$};
				\clip (\xmin+0.1,\ymin+0.1) rectangle (\xmax-0.1,\ymax-0.1);
				\draw (-1,-1)--(2,3);
				\draw[dashed] (-1,0)node[above]{$-1$}--(-1,-1)node[below left]{$A$}--(0,-1)node[right]{$-1$} 
				(2,0)node[below]{$2$}--(2,3)node[above right]{$B$}--(0,3)node[left]{$3$};
				\fill (0,0) circle (1.0pt) (-1,-1) circle (1.0pt) (2,3) circle (1.0pt);
			\end{tikzpicture}
		\end{center}		
		Ta có $MA+MB\geq AB$.\\
		Dấu \lq\lq=\rq\rq\, xảy ra khi $M\equiv I$ là giao điểm của $AB$ với trục $Ox$.\\
		Ta có: $\overrightarrow{AI}=(t+1;1),\overrightarrow{AB}=(3;4)$ mà $\overrightarrow{AI},\overrightarrow{AB}$ cùng phương nên 
		$\dfrac{t+1}{3}=\dfrac{1}{4}\Rightarrow t=-\dfrac{1}{4}\Leftrightarrow\sin x=\dfrac{-1}{4}$.
		Do đó $a=-1,b=4\Rightarrow a^3+b^2=(-1)^3+4^2=15.$}
\end{ex}

\begin{ex}%[Thi thử TN - Chuyên Hạ Long-Quảng Ninh-L2-23]%[Duong Xuan Loi-EX6]%[2D2B5-1]
	Phương trình $\left(x^2-4x+3\right)\log_{2023}\left(x^2-4\right)=0$ có bao nhiêu nghiệm?
	\choice
	{$0$}
	{$2$}
	{\True $3$}
	{$4$}
	\loigiai{
		Điều kiện: $x^2-4>0\Leftrightarrow\hoac{&x>2\\&x <-2.}$ \\
		Phương trình tương đương $\hoac{&x^2-4x+3=0\\&\log_{2023}\left(x^2-4\right)=0}\Leftrightarrow\hoac{&x=1\\&x=3\\&x=\sqrt{5}\\&x=-\sqrt{5}.}$ \\
		Đối chiếu điều kiện, ta nhận $x=3,x=\pm\sqrt{5}$ nên phương trình đã cho có ba nghiệm.}
\end{ex}

\begin{ex}%[Thi thử TN - Chuyên Hạ Long-Quảng Ninh-L2-23]%[Duong Xuan Loi-EX6]%[2D2B5-2]
	Tổng các nghiệm của phương trình $3\cdot 9^x-5\cdot 6^x+2\cdot 4^x=0$ là
	\choice
	{$0$}
	{\True $-1$}
	{$\dfrac{5}{3}$}
	{$1$}
	\loigiai{
		Chia cả hai vế cho $4^x$ ta được\\
		$3\cdot 9^x-5\cdot 6^x+2\cdot 4^x=0\Leftrightarrow 3\cdot\left(\dfrac{3}{2}\right)^{2x}-5\cdot\left(\dfrac{3}{2}\right)^x+2=0\Leftrightarrow\hoac{&\left(\dfrac{3}{2}\right)^x=1\\&\left(\dfrac{3}{2}\right)^x=\dfrac{2}{3}}\Leftrightarrow\hoac{&x=0\\&x=-1.}$}
\end{ex}

\begin{ex}%[Thi thử TN - Chuyên Hạ Long-Quảng Ninh-L2-23]%[Duong Xuan Loi-EX6]%[2D2B6-2]
	Bất phương trình $25^x-6\cdot 5^x+5\leq 0$ có tập nghiệm là $[a;b]$. Tính $a\cdot b$. 
	\choice
	{$1$}
	{\True $0$}
	{$6$}
	{$5$}
	\loigiai{
		$25^x-6\cdot 5^x+5\leq 0\Leftrightarrow(5^x)^2-6\cdot 5^x+5\leq 0$.\\
		Đặt $t=5^x$, $t>0$.\\
		Bất phương trình có dạng $t^2-6t+5\leq 0\Leftrightarrow 1\leq t\leq 5\Rightarrow 1\leq 5^x\leq 5\Leftrightarrow 0\leq x\leq 1$.\\
		Tập nghiệm bất phương trình là $[0;1]\Rightarrow \heva{&a=0\\&b=1}\Rightarrow a\cdot b=0$.
	}
\end{ex}

\begin{ex}%[Thi thử TN - Chuyên Hạ Long-Quảng Ninh-L2-23]%[Duong Xuan Loi-EX6]%[2D3B1-3]
	Chọn khẳng định \textbf{sai}? 
	\choice
	{$\displaystyle\int x^2\ln(x-2)\mathrm{\,d}x=\dfrac{x^3}{3}\ln(x-2)+\dfrac{1}{3}\displaystyle\int\dfrac{x^3}{2-x}\mathrm{\,d}x,\forall x\in(2;+\infty)$}
	{\True $\displaystyle\int x^2\ln(x-2)\mathrm{\,d}x=\dfrac{x^3}{3}\ln(x-2)+\dfrac{1}{3}\displaystyle\int\dfrac{x^2}{x-2}\mathrm{\,d}x,\forall x\in(2;+\infty)$}
	{$\displaystyle\int x^2\ln(x-2)\mathrm{\,d}x=\dfrac{x^3-8}{3}\ln(x-2)-\displaystyle\int\dfrac{x^2+2x+4}{3}\mathrm{\,d}x,\forall x\in(2;+\infty)$}
	{$\displaystyle\int x^2\ln(x-2)\mathrm{\,d}x=\dfrac{x^3}{3}\ln(x-2)-\dfrac{1}{3}\displaystyle\int\dfrac{x^3}{x-2}\mathrm{\,d}x,\forall x\in(2;+\infty)$}
	\loigiai{
		Đặt $\heva{&u=\ln(x-2)\\&\mathrm{\,d}v=x^2\mathrm{\,d}x}\Rightarrow\heva{&\mathrm{\,d}u=\dfrac{1}{x-2}\mathrm{\,d}x\\&v=\dfrac{x^3}{3}.}$ \\
		Suy ra $\displaystyle\int x^2\ln(x-2)\mathrm{\,d}x=\dfrac{x^3}{3}\ln(x-2)-\dfrac{1}{3}\displaystyle\int\dfrac{x^3}{x-2}\mathrm{\,d}x,\forall x\in(2;+\infty)$.}
\end{ex}

\begin{ex}%[Thi thử TN - Chuyên Hạ Long-Quảng Ninh-L2-23]%[Duong Xuan Loi-EX6]%[2D3B2-2]
	Biết $\displaystyle\int\limits_0^1 x\mathrm{e}^{3x^2+3}\mathrm{\,d}x=\dfrac{\mathrm{e}^3}{a}\left(\mathrm{e}^b-\mathrm{e}^c\right)\left(a,b,c\in\mathbb{Z}\right)$. Tính $a+b+c$. 
	\choice
	{$6$}
	{$0$}
	{\True $9$}
	{$3$}
	\loigiai{
		Đặt $t=3x^2+3\Rightarrow\mathrm{\,d}t=6x\mathrm{\,d}x\Rightarrow x\mathrm{\,d}x=\dfrac{\mathrm{\,d}t}{6}$.\\
		Đổi cận 
		
		\begin{tabular}{|c|c|c|}
			\hline
			$x$&  $0$ & $1$\\
			\hline
			$t$& $3$ & $6$\\
			\hline
		\end{tabular}\\
		$I=\displaystyle\int\limits_3^6\mathrm{e}^t\cdot\dfrac{\mathrm{\,d}t}{6}=\dfrac{1}{6}\cdot\mathrm{e}^t\bigg|_3^6 =\dfrac{1}{6}\cdot\left(\mathrm{e}^6-\mathrm{e}^3\right)=\dfrac{\mathrm{e}^3}{6}\left(\mathrm{e}^3-1\right)$.\\
		$\begin{aligned}&\Rightarrow a=6; b=3;c=0\\&\Rightarrow a+b+c=6+3+0=9.\end{aligned}$}
\end{ex}

\begin{ex}%[Thi thử TN - Chuyên Hạ Long-Quảng Ninh-L2-23]%[Duong Xuan Loi-EX6]%[2D3B2-2]
	Cho hàm số bậc nhất $f(x)$ thỏa mãn $\displaystyle\int\limits_0^1 f(x)\mathrm{\,d}x=4$; $\displaystyle\int\limits_2^3 f(x)\mathrm{\,d}x=2$. Tính $I=\displaystyle\int\limits_0^1 f\left(f(2x-5)\right)\mathrm{\,d}x$. 
	\choice
	{$6$}
	{$\dfrac{7}{2}$}
	{\True $-4$}
	{$\dfrac{3}{2}$}
	\loigiai{
		Hàm số bậc nhất $f(x)=ax+b$.\\
		$4=\displaystyle\int\limits_0^1 f(x)\mathrm{\,d}x=\displaystyle\int\limits_0^1(ax+b)\mathrm{\,d}x=\left(\dfrac{ax^2}{2}+bx\right)\bigg|_0^1=\dfrac{a}{2}+b\Rightarrow\dfrac{a}{2}+b=4.\quad(1)$\\
		$2=\displaystyle\int\limits_2^3(ax+b)\mathrm{\,d}x=\left(\dfrac{ax^2}{2}+bx\right)\bigg|_2^3=\dfrac{9a}{2}+3b-\dfrac{4a}{2}-2b=\dfrac{5a}{2}+b\Rightarrow\dfrac{5a}{2}+b=2.\quad(2)$\\
		Từ $(1)$ và $(2)$ ta có $\heva{&\dfrac{a}{2}+b=4\\&\dfrac{5a}{2}+b=2}\Leftrightarrow\heva{&a=-1\\&b=\dfrac{9}{2}}\Rightarrow f(x)=-x+\dfrac{9}{2.}$ \\
		$f(2x-5)=-2x+5+\dfrac{9}{2}$.\\
		$f\left(f(2x-5)\right)=2x-5-\dfrac{9}{2}+\dfrac{9}{2}=2x-5$.\\
		$I=\displaystyle\int\limits_0^1 f\left(f(2x-5)\right)\mathrm{\,d}x=\displaystyle\int\limits_0^1(2x-5)\mathrm{\,d}x=-4$.}
\end{ex}

\begin{ex}%[Thi thử TN - Chuyên Hạ Long-Quảng Ninh-L2-23]%[Duong Xuan Loi-EX6]%[2D4B2-4]
	Tập hợp các điểm biểu diễn số phức $z$ thỏa mãn $|z+4|+\left|\overline{z}-4\right|=8$ trên mặt phẳng phức là 
	\choice
	{Đường Elip}
	{Đường thẳng}
	{\True Đoạn thẳng}
	{Tập rỗng}
	\loigiai{
		Gọi $M(x;y)$ là điểm biểu diễn của số phức $z=x+yi(x,y\in\mathbb{R})$ và $\overline{z}=x-yi$.\\
		Ta có: $|z+4|+\left|\overline{z}-4\right|=8\Leftrightarrow|x+4+yi|+|x-4-yi|=8$ \\
		$ \Leftrightarrow\sqrt{(x+4)^2+y^2}+\sqrt{(x-4)^2+y^2}=8 $.\\
		Đặt $A(-4;0),B(4;0)\Rightarrow AB=\sqrt{(4+4)^2+0^2}=8$.\\
		Nên $MA+MB=AB$.\\
		Do đó tập hợp các điểm biểu diễn của $z$ chính là đoạn thẳng $AB$.}
\end{ex}

\begin{ex}%[Thi thử TN - Chuyên Hạ Long-Quảng Ninh-L2-23]%[Duong Xuan Loi-EX6]%[2H1B3-2]
	Cho hình chóp $S.ABC$ đáy tam giác $ABC$ vuông ở $A$, $\widehat{SCA}=\widehat{SBA}=90^{\circ}$, $SA=13a$; $AB=3a$; $BC=5a$. Tính thể tích chóp $S.ABC$. 
	\choice
	{\True $24a^3$}
	{$48a^3$}
	{$72a^3$}
	{$26a^3$}
	\loigiai{
		\immini{
			Gọi $H$ là hình chiếu của $S$ trên $(ABC)$.\\
			Có $\heva{&AB\perp SB\\&AB\perp SH}\Rightarrow AB\perp(SHB)\Rightarrow AB\perp HB$.\\
			Mặt khác trong $(ABC)$, ta lại có  $AB\perp AC\Rightarrow AC\parallel HB$.\\
			$\heva{&AC\perp SC\\&AC\perp SH}\Rightarrow AC\perp(SHC)\Rightarrow AC\perp HC.$\\
			Trong $(ABC)$, có $AB\perp AC\Rightarrow  AB\parallel HC$.\\
			Suy ra tứ giác $ABHC$ là hình chữ nhật nên $AH=BC=5a$.\\
			$AC=\sqrt{BC^2-AB^2}=4a$.\\
			Tam giác SHA vuông tại $H$ có\\
			$SH=\sqrt{SA^2-AH^2}=12a$.\\
			Vậy $V_{S.ABC}=\dfrac{1}{3}SH\cdot\dfrac{1}{2}AB\cdot AC=24a^3$.
		}{
			\begin{tikzpicture}[scale=1, font=\footnotesize, line join=round, line cap=round,>=stealth]
				\path
				(0,0) coordinate (A)
				(-1,-1.6) coordinate (C)
				(4,0)coordinate (B)
				($(B)+(C)-(A)$) coordinate (H)
				($(H)+(0,4.5)$) coordinate (S)
				;
				\draw (S)--(B)--(H)--(C)--(A)--cycle (C)--(S)--(H);
				\draw[dashed] (A)--(B)--(C);	
				\foreach \p/\q in {S/90,A/180,B/-90,C/-90,H/0}			
				\fill[black] (\p) circle (1.0pt)node[shift={(\q:2.5mm)}]{$\p$};
				\draw pic[draw=black,angle radius=0.2cm] {right angle = S--B--A}
				pic[draw=black,angle radius=0.2cm] {right angle = B--A--C}
				pic[draw=black,angle radius=0.2cm] {right angle = S--H--C}
				pic[draw=black,angle radius=0.2cm] {right angle = S--H--B}; 
			\end{tikzpicture}
	}}
\end{ex}

\begin{ex}%[Thi thử TN - Chuyên Hạ Long-Quảng Ninh-L2-23]%[Duong Xuan Loi-EX6]%[2H3B1-1]
	Trong không gian cho hệ trục $Oxyz$; cho $A(1;1;2)$, $B(-4;0;11)$, $C(0;-21;0)$. Có bao nhiêu điểm $D$ sao cho $A$, $B$, $C$, $D$ là bốn đỉnh của một hình bình hành?
	\choice
	{Có vô số điểm $D$}
	{Có $3$ điểm $D$}
	{\True Có duy nhất một điểm $D$}
	{Có $2$ điểm $D$}
	\loigiai{
		Ta có: $\overrightarrow{AB}=(-5;-1;9)$; $\overrightarrow{AC}=(-1;-22;-2)\Rightarrow A,B,C$ không thẳng hàng.\\
		Có duy nhất một điểm $D$ để $A, B, C, D$ là bốn đỉnh của một hình bình hành.}
\end{ex}

\begin{ex}%[Thi thử TN - Chuyên Hạ Long-Quảng Ninh-L2-23]%[Duong Xuan Loi-EX6]%[2H3B1-3]
	Trong không gian cho hệ trục $Oxyz$, phương trình mặt cầu tâm $O$, tiếp xúc với mặt phẳng $(P)\colon 2x+y+2z-9=0$ là 
	\choice
	{$x^2+y^2+x^2=6$}
	{$x^2+y^2+x^2=81$}
	{$x^2+y^2+x^2=3$}
	{\True $x^2+y^2+x^2=9$}
	\loigiai{
		Ta có $R=\mathrm{d}\left(O;(P)\right)=\dfrac{|-9|}{\sqrt{2^2+1^2+2^2}}=3$.\\
		Phương trình mặt cầu tâm $O$, tiếp xúc với mặt phẳng $(P)\colon 2x+y+2z-9=0$ là 
		$x^2+y^2+x^2=9$.}
\end{ex}

\begin{ex}%[Thi thử TN - Chuyên Hạ Long-Quảng Ninh-L2-23]%[Duong Xuan Loi-EX6]%[2H3B2-7]
	Trong không gian cho hệ trục $Oxyz$; mặt phẳng $(P)$ có phương trình $\dfrac{x}{2}+\dfrac{y}{3}-\dfrac{z}{6}-1=0$ cắt trục $Oy$ tại điểm $A$ có tọa độ 
	\choice
	{\True $A(0;3;0)$}
	{$A\left(0;\dfrac{1}{3};0\right)$}
	{$A\left(\dfrac{1}{2};\dfrac{1}{3};\dfrac{-1}{6}\right)$}
	{$A\left(\dfrac{1}{2};0;\dfrac{-1}{6}\right)$}
	\loigiai{
		Gọi $Oy\cap(P)=A(0;y;0)$.\\
		Xét phương trình: $\dfrac{0}{2}+\dfrac{y}{3}-\dfrac{0}{6}-1=0\Rightarrow y=3$.\\
		Vậy $A(0;3;0)$.}
\end{ex}

\begin{ex}%[Thi thử TN - Chuyên Hạ Long-Quảng Ninh-L2-23]%[Duong Xuan Loi-EX6]%[2H3B3-4]
	Trong không gian có hệ trục $Oxyz$, đường thẳng $(d)$ có phương trình: $\heva{&x=2-t\\&y=3+t\\&z=1+t}$ và mặt phẳng $(\Delta)\colon x+2y-2z+3=0$. Góc giữa $(d)$ và $(\Delta)$ bằng $\alpha$. Khẳng định nào sau đây đúng?
	\choice
	{$\cot\alpha=\dfrac{\sqrt{3}}{9}$}
	{$\tan\alpha=\dfrac{\sqrt{3}}{9}$}
	{\True $\sin\alpha=\dfrac{\sqrt{3}}{9}$}
	{$\cos\alpha=\dfrac{\sqrt{3}}{9}$}
	\loigiai{
		Đường thẳng $(d)$ có véc-tơ chỉ phương $\overrightarrow{u}=(-1;1;1)$, mặt phẳng $(\Delta)$ có véc-tơ pháp tuyến $\overrightarrow{n}=(1;2;-2)$, khi đó:\\
		$\sin\left((d),(\Delta)\right)=\sin\alpha=\dfrac{\left|\overrightarrow{u}\cdot\overrightarrow{n}\right|}{\left|\overrightarrow{u}\right|\cdot\left|\overrightarrow{n}\right|}=|\cos (\overrightarrow{u},\overrightarrow{n})|=\dfrac{\left|(-1)\cdot 1+1\cdot 2+1\cdot (-2)\right|}{\sqrt{(-1)^2+1^2+1^2}\cdot\sqrt{1^2+2^2+(-2)^2}}=\dfrac{\sqrt{3}}{9}$.}
\end{ex}

\begin{ex}%[Thi thử TN - Chuyên Hạ Long-Quảng Ninh-L2-23]%[Duong Xuan Loi-EX6]%[2H3B3-6]
	Trong không gian cho hệ trục $Oxyz$, cho đường thẳng $\Delta\colon\dfrac{x-1}{2}=\dfrac{y+1}{1}=\dfrac{z-3}{2}$ và mặt phẳng $(\alpha)\colon 2x+y+3z-10=0$. Gọi $M(a;b;c)$ là giao điểm của $(\alpha)$ và $\Delta$. Tính $2a+b+c$. 
	\choice
	{$5$}
	{\True $4$}
	{$2$}
	{$3$}
	\loigiai{
		Vì $M\in\Delta\Rightarrow M\left(1+2t;-1+t;3+2t\right)$.\\
		$M\in(\alpha)\colon 2(1+2t)-1+t+3(3+2t)-10=0\Leftrightarrow 11t=0\Leftrightarrow t=0$. \\
		Suy ra $M(1;-1;3)\Rightarrow 2a+b+c=2-1+3=4 $.}
\end{ex}

\begin{ex}%[Thi thử TN - Chuyên Hạ Long-Quảng Ninh-L2-23]%[Duong Xuan Loi-EX6]%[2D1K2-4]
	Có bao nhiêu số nguyên $m\in[-2023;2023]$ để đồ thị hàm số $y=x^3+mx^2-m(2m+1)x+m^2$ có hai điểm cực trị nằm về hai phía của trục hoành?
	\choice
	{$4044$}
	{\True $4045$}
	{$4046$}
	{$4047$}
	\loigiai{
		Đồ thị hàm số đã có hai điểm cực trị nằm về hai phía của trục hoành khi và chỉ khi phương trình $y=0$ có ba nghiệm phân biệt.\\
		Ta có $y=0\Leftrightarrow(x-m)\left(x^2+2mx-m\right)=0\Leftrightarrow\hoac{&x=m\\&g(x)=x^2+2mx-m=0.}$ \\
		Yêu cầu bài toán tương đương $g(x)=0$ có hai nghiệm phân biệt khác $m$ tương đương
		$$\heva{&\Delta'_g=m^2+m>0\\&g(m)\neq 0}\Leftrightarrow\heva{&\hoac{&m>0\\&m <-1}\\&m\neq 0; m\neq\dfrac{1}{3}}\Leftrightarrow\heva{&\hoac{&m>0\\&m <-1}\\& m\neq\dfrac{1}{3}.}$$
		Vì $m$ nguyên và $m\in[-2023;2023]$ nên có $4045$ giá trị.}
\end{ex}

\begin{ex}%[Thi thử TN - Chuyên Hạ Long-Quảng Ninh-L2-23]%[Duong Xuan Loi-EX6]%[2D1K5-3]
	Cho hàm số $y=f(x)$ bậc ba có bảng biến thiên sau 
	\begin{center}
		\begin{tikzpicture}
			\tkzTabInit[nocadre=false,lgt=1.2,espcl=2.5,deltacl=0.6]
			{$x$ /1, $y'$ /1, $y$ /2.5}
			{$-\infty$,$0$,$\dfrac{3}{2}$,$+\infty$}
			\tkzTabLine{,+,$0$,-,$0$,+,}
			\tkzTabVar{-/$-\infty$,+/$1$,-/$0$,+/$+\infty$}
		\end{tikzpicture}
	\end{center}
	Đặt $g(x)=f\left(\dfrac{1}{2}f^2(x)-f(x)\right)$. Phương trình $g'(x)=0$ có bao nhiêu nghiệm phân biệt?
	\choice
	{$8$}
	{$6$}
	{$9$}
	{\True $7$}
	\loigiai{
		Ta có $g'(x)=f'(x)[f(x)-1]f'\left(\dfrac{1}{2}f^2(x)-f(x)\right)$ nên 
		$g'(x)=0\Leftrightarrow\hoac{&f'(x)=0; f(x)=1\\&f'\left(\dfrac{1}{2}f^2(x)-f(x)\right)=0.}$ \\
		Xét phương trình $f'(x)=0\Leftrightarrow\hoac{&x=0\\&x=\dfrac{3}{2}.}$ \\
		Xét phương trình $f(x)=1\Leftrightarrow\hoac{&x=0\\&x=a\left(a>\dfrac{3}{2}\right).}$ \\
		Xét phương trình\\
		$f'\left(\dfrac{1}{2}f^2(x)-f(x)\right)=0\cdot\Leftrightarrow\hoac{&\dfrac{1}{2}f^2(x)-f(x)=0\\&\dfrac{1}{2}f^2(x)-f(x)=\dfrac{3}{2}}\Leftrightarrow\hoac{&f(x)=0\Leftrightarrow\hoac{&x=\dfrac{3}{2}\\&x=b(b<0)}\\&f(x)=2\Leftrightarrow x=c(c>a)\\&f(x)=3\Leftrightarrow x=\mathrm{d}(d>c)\\&f(x)=-1\Leftrightarrow x=e(e<b).}$ \\
		Vậy phương trình $g'(x)=0$ có $7$ nghiệm phân biệt.}
\end{ex}

\begin{ex}%[Thi thử TN - Chuyên Hạ Long-Quảng Ninh-L2-23]%[Duong Xuan Loi-EX6]%[2D2K6-2]
	Tổng tất cả các giá trị nguyên của $m$ thỏa $\log_3\left|\dfrac{x^2-2(m-1)x+2-m}{3x^2+2x+1}\right|\leq 1$ với mọi $x$ là
	\choice
	{$-14$}
	{$21$}
	{$14$}
	{\True Đáp án khác}
	\loigiai{
		Ta có $\log_3\left|\dfrac{x^2-2(m-1)x+2-m}{3x^2+2x+1}\right|\leq 1$ với mọi $x$ \\
		$ \Leftrightarrow\heva{&x^2-2(m-1)x+2-m\neq 0\\&\left|x^2-2(m-1)x+2-m\right|\leq 3\left(3x^2+2x+1\right)} $ với mọi $x$ \\
		$ \Leftrightarrow\heva{&x^2-2(m-1)x+2-m\neq 0\\&x^2-2(m-1)x+2-m\leq 9x^2+6x+3\\&x^2-2(m-1)x+2-m\geq-9x^2-6x-3} $ với mọi $x$ \\
		$ \Leftrightarrow\heva{&x^2-2(m-1)x+2-m\neq 0\\&-8x^2-2(m+2)x-1-m\leq 0\\&10x^2-2(m-4)x+5-m\geq 0} $ với mọi $x$ \\
		$ \Leftrightarrow\heva{&(m-1)^2+m-2<0\\&(m+2)^2-8(m+1)\leq 0\\&(m-4)^2-10(5-m)\leq 0}\Leftrightarrow\heva{&m^2-m-1<0\\&(m+2)^2-8(m+1)\leq 0\\&(m-4)^2-10(5-m)\leq 0.} $ \\
		Mà $m$ là số nguyên nên $m=0$ hay $m=1$.\\
		Vậy tổng tất cả các giá trị $m$ là $1$.}
\end{ex}

\begin{ex}%[Thi thử TN - Chuyên Hạ Long-Quảng Ninh-L2-23]%[Duong Xuan Loi-EX6]%[2D3K2-4]
	Cho hàm số $f(x)$ có đạo hàm trên $\mathbb{R}$ thỏa mãn $f(x)+f'(x)=x^3+3x^2-4x+4,\forall x\in\mathbb{R}$ và $f(1)=5$. Tính diện tích hình phẳng giới hạn bởi đồ thị hàm số $y=f(x)$ và $y=f'(x)$. 
	\choice
	{\True $\dfrac{131}{4}$}
	{$\dfrac{125}{4}$}
	{$\dfrac{35}{4}$}
	{$\dfrac{203}{4}$}
	\loigiai{
		Chọn $f(x)=ax^3+bx^2+cx+d\Rightarrow f'(x)=3ax^2+2bx+c$.\\
		Suy ra $f(x)+f'(x)=ax^3+(3a+b)x^2+(2b+c)x+d+c$.\\
		Từ giả thiết ta có: $\heva{&a=1\\&3a+b=3\\&2b+c=-4\\&d+c=4}\Rightarrow\heva{&a=1\\&b=0\\&c=-4\\&d=8}\Rightarrow\heva{&f(x)=x^3-4x+8\\&f'(x)=3x^2-4}$ (thỏa mãn).\\
		Phương trình hoành độ giao điểm của hàm số $y=f(x),y=f'(x)$ là\\
		$f(x)=f'(x)\Leftrightarrow x^3-4x+8=3x^2-4\Leftrightarrow x^3-3x^2-4x+12=0\Leftrightarrow\hoac{&x=-2\\&x=2\\&x=3.}$ \\
		Khi đó, diện tích hình phẳng giới hạn bởi đồ thị hàm số $y=f(x)$ và $y=f'(x)$ bằng\\
		$S=\displaystyle\int\limits_{-2}^3\left|x^3-3x^2-4x+12\right|\mathrm{\,d}x=\displaystyle\int\limits_{-2}^2\left(x^3-3x^2-4x+12\right)\mathrm{\,d}x+\displaystyle\int\limits_2^3\left(-x^3+3x^2+4x-12\right)=\dfrac{131}{4}$.}
\end{ex}

\begin{ex}%[Thi thử TN - Chuyên Hạ Long-Quảng Ninh-L2-23]%[Duong Xuan Loi-EX6]%[2D4K3-1]
	Cho $z+\dfrac{1}{z}=-1$. Tính $P=\left|z^{2023}+\dfrac{1}{z^{2023}}\right|$. 
	\choice
	{$P=\sqrt{2}$}
	{\True $P=1$}
	{$P=0$}
	{$P=-1$}
	\loigiai{
		$z+\dfrac{1}{z}=-1\Leftrightarrow z^2+z+1=0\Leftrightarrow(z-1)\left(z^2+z+1\right)=0\Leftrightarrow z^3-1=0\Leftrightarrow z^3=1$.\\
		Khi đó $P=\left|z^{2023}+\dfrac{1}{z^{2023}}\right|=\left|(z^3)^{674}z+\dfrac{1}{(z^3)^{674}z}\right|=\left|z+\dfrac{1}{z}\right|=1$.}
\end{ex}

\begin{ex}%[Thi thử TN - Chuyên Hạ Long-Quảng Ninh-L2-23]%[Duong Xuan Loi-EX6]%[2D4K5-1]
	Cho hai số phức $z_1$; $z_2$ thỏa mãn $|z-2-i|=5$; $|z+2+mi|=|z-m+i|,(m\in\mathbb{R})$. Giá trị nhỏ nhất của $P=|z_1-z_2|$ thuộc đoạn nào sau đây?
	\choice
	{$[4;5]$}
	{\True $[8;9]$}
	{$[5;6]$}
	{$[6;7]$}
	\loigiai{
		\begin{center}
			\begin{tikzpicture}[scale=1,font=\footnotesize,line join=round,line cap=round,>=stealth]
				\path
				(0,0) coordinate (I)
				($(I)+(-160:2.5cm)$) coordinate (A)
				($(I)+(-60:2.5cm)$) coordinate (B)
				($(A)!0.5!(B)$) coordinate (K)			
				;
				\draw (K)--(I) circle(2.5cm) ($(A)!1.2!(B)$)node[right]{$d$}--($(B)!1.2!(A)$) (I)--(B);
				\foreach \p/\q in {I/90,A/230,B/-85,K/-100}
				\fill[black] (\p) circle (1.0pt) ($(\p)+(\q:2.5mm)$) node{$\p$};		
				\draw pic[draw=black,angle radius=0.2cm] {right angle = I--K--B}; 
			\end{tikzpicture}
		\end{center}		
		Đặt $z=x+yi,(x,y\in\mathbb{R})$. Khi đó:\\
		$\begin{aligned}&|z-2-i|=5\Leftrightarrow(x-2)^2+(y-1)^2=25\\&|z+2+mi|=|z-m+i|\Leftrightarrow(2m+4)x+(2m-2)y+3=0\end{aligned}$ \\
		$ \Rightarrow z_1, z_2 $ vừa thuộc đường tròn $(C)$ có tâm $I(2;1)$ và bán kính $R=5$ vừa thuộc đường thẳng $d$ có phương trình $(2m+4)x+(2m-2)y+3=0$ (ở đây, $d$ không đi qua tâm $I$ mà luôn đi qua điểm $K\left(-\dfrac{1}{2};\dfrac{1}{2}\right)$ cố định) và giả sử $d$ cắt $(C)$ tại hai điểm $A, B$.\\
		Do đó, $P=|z_1-z_2|=AB$. Vậy $P$ có giá trị nhỏ nhất khi $d\perp IK$. Khi đó,\\
		$AB=2KB=2\sqrt{R^2-IK^2}=\sqrt{74}\Rightarrow$ giá trị nhỏ nhất của $P=|z_1-z_2|$ thuộc đoạn $[8;9]$.}
\end{ex}

\begin{ex}%[Thi thử TN - Chuyên Hạ Long-Quảng Ninh-L2-23]%[Duong Xuan Loi-EX6]%[2H1K3-2]
	Cho hình chóp $S.ABC$ có đáy $ABC$ là tam giác cân ở $B$; $\widehat{ABC}=120^{\circ}$; $AC=a\sqrt{3}$. Các cạnh bên $SA=SB=SC$; $SB$ tạo với mặt đáy một góc $60^{\circ}$. Tính thể tích khối chóp $S.ABC$. 
	\choice
	{$\dfrac{a^3\sqrt{3}}{4}$}
	{\True $\dfrac{a^3}{4}$}
	{$\dfrac{3a^3}{4}$}
	{$a^3\sqrt{3}$}
	\loigiai{		
		\immini{
			$SA=SB=SC\Rightarrow$ chân đường cao $H$ kẻ từ đỉnh $S$ của hình chóp $S.ABC$ chính là tâm đường tròn ngoại tiếp tam giác $ABC$ (với $H\in BM, M$ là trung điểm của $AC$).\\
			Khi đó, $SH\perp(ABC)\Rightarrow\left(SB,(ABC)\right)=\left(SB,HB\right)=\widehat{SBH}=60^{\circ}$.\\
			Ta có $AC^2=AB^2+BC^2-2AB\cdot BC\cdot\cos\widehat{ABC}\Rightarrow AB=BC=a$ \\
			$\Rightarrow S_{\triangle ABC}=\dfrac{1}{2}\cdot AB\cdot BC\cdot\sin\widehat{ABC}=\dfrac{1}{2}\cdot a\cdot a\cdot\sin 120^{\circ}=\dfrac{a^2\sqrt{3}}{4}$ \\
			$\Rightarrow BH=\dfrac{AC}{2\sin\widehat{ABC}}=a\Rightarrow SH=BH\cdot\tan\widehat{SBH}=a\sqrt{3}$.
		}{
			\begin{tikzpicture}[scale=1, font=\footnotesize, line join=round, line cap=round,>=stealth]
				\path
				(0,0) coordinate (B)
				(1.2,-1.6) coordinate (A)
				(5,0) coordinate (C)
				($(A)!0.5!(C)$) coordinate (M)
				($(B)!2/3!(M)$) coordinate (H)
				($(H)+(0,3.3)$) coordinate (S)
				;
				\draw (S)--(B)--(A)--(C)--cycle (S)--(A);	
				\draw[dashed] (C)--(B)--(M) (S)--(H);	
				\foreach \p/\q in {S/90,B/180,A/-90,C/0,H/-90,M/-40}			
				\fill[black] (\p) node[shift={(\q:3mm)}]{$\p$} circle (1.0pt);
				\draw pic[draw=black,angle radius=0.2cm] {right angle = S--H--M}
				pic[draw=black,angle radius=0.2cm] {right angle = B--M--C}
				; 
			\end{tikzpicture}
		}\noindent
		Vậy $V_{S.ABC}=\dfrac{1}{3}\cdot S_{\triangle ABC}\cdot SH=\dfrac{a^3}{4}$.}
\end{ex}

\begin{ex}%[Thi thử TN - Chuyên Hạ Long-Quảng Ninh-L2-23]%[Duong Xuan Loi-EX6]%[2H3K3-3]
	Trong không gian cho hệ trục $Oxyz$, cho đường thẳng $d\colon\heva{&x=1+t\\&y=2-t\\&z=-1-2t}$ và mặt phẳng $(P)\colon x-2y+z-1=0$. Đường thẳng $\Delta$ nằm trong mặt phẳng $(P)$; cắt và vuông góc với đường thẳng $d$. Đường thẳng $\Delta$ không đi qua điểm nào dưới đây. 
	\choice
	{$E(6;-3;-11)$}
	{$F(11;0;-10)$}
	{$G(1;-6;-12)$}
	{\True $A(2;7;13)$}
	\loigiai{
		Xét phương trình $1+t-2(2-t)-1-2t-1=0\Leftrightarrow t=5\Rightarrow d\cap(P)=M(6;-3;-11)$.\\
		$d$ có véc-tơ chỉ phương $\overrightarrow{u}=(1;-1;-2)$.\\
		Mặt phẳng $(P)$ có véc-tơ pháp tuyến $\overrightarrow{n}=(1;-2;1)$.\\
		Do $\heva{&\Delta\subset(P)\\&\Delta\perp d,\Delta\cap d=M}\Rightarrow\left[\overrightarrow{u},\overrightarrow{n}\right]=(-5;-3;-1)$ là véc-tơ chỉ phương của $\Delta$.\\
		PTĐT $\Delta\colon\dfrac{x-6}{5}=\dfrac{y+3}{3}=\dfrac{z+11}{1}$.\\
		Thay tọa độ điểm $A(2;7;13)$ vào phương trình của $\Delta$ ta có $\Delta\colon\dfrac{2-6}{5}=\dfrac{7+3}{3}=\dfrac{13+11}{1}\Rightarrow\dfrac{-4}{5}=\dfrac{10}{3}$.\\
		Vô lý. Do đó $A\notin\Delta$.}
\end{ex}

\begin{ex}%[Thi thử TN - Chuyên Hạ Long-Quảng Ninh-L2-23]%[Duong Xuan Loi-EX6]%[2H3K3-7]
	Cho mặt cầu $S(O;9)$. Một hình nón có đỉnh và đường tròn đáy nằm trên mặt cầu $(S)$. Khi thể tích của hình nón lớn nhất, diện tích đường tròn đáy của hình nón thuộc khoảng nào dưới đấy?
	\choice
	{$[200;220]$}
	{$(230;240)$}
	{$(200;220)$}
	{\True $[220;230]$}
	\loigiai{
		\begin{center}
			\begin{tikzpicture}[scale=0.7, font=\footnotesize, line join=round, line cap=round, >=stealth]
				\def\vitri{-30} 
				\path
				(0,0) coordinate (O)
				(\vitri:3) coordinate (A)
				(-180-\vitri:3) coordinate (B)
				($(A)!0.5!(B)$) coordinate (H)				
				(90:3) coordinate (S)
				;				
				\draw (O) circle (3) (A)arc (0:-180: {3*cos(\vitri)} and {0.5*cos(\vitri)});
				\draw[dashed] (H)--(S)--(A)arc (0:180: {3*cos(\vitri)} and {0.5*cos(\vitri)})--(B)--(S) (A)--(B);
				\foreach \p/\q in {O/0,A/-10,B/190,S/90}
				\fill[black] (\p) circle (1.0pt) ($(\p)+(\q:2.5mm)$) node{$\p$};
				\draw pic[draw=black,angle radius=0.2cm] {right angle = S--H--A}; 
			\end{tikzpicture}
		\end{center}
		Gọi khoảng cách từ tâm của khối cầu đến đáy của hình nón là $x$. $(0<x<9)$.\\
		Khi đó, chiều cao của khối nón $h=x+9$ và bán kính của khối nón $r=\sqrt{9^2-x^2}\Rightarrow r^2=81-x^2$.\\
		Thể tích khối nón là $V=\dfrac{1}{3}S\cdot h=\dfrac{1}{3}\pi\left(81-x^2\right)(x+9)=\dfrac{1}{3}\pi(x+9)^2(9-x)$.\\
		Xét hàm số $f(x)=\dfrac{1}{3}\pi(x+9)^2\cdot (9-x)\Rightarrow f'(x)=\dfrac{1}{3}\pi\left(-3x^2-18x+81\right)=0\Leftrightarrow x=3$.\\
		Bảng biến thiên: 
		\begin{center}
			\begin{tikzpicture}
				\tkzTabInit[nocadre=false,lgt=1.2,espcl=3,deltacl=0.6]
				{$x$ /0.6, $f'(x)$ /0.6, $f(x)$ /2.1}
				{$0$,$3$,$9$}
				\tkzTabLine{,+,$0$,-,}
				\tkzTabVar{-/, +/,-/}
			\end{tikzpicture}
		\end{center}
		Vậy $V_{\max}\Leftrightarrow x=3$. Khi đó diện tích đáy là $S=\pi\left(9^2-3^2\right)=72\pi\in[220;230]$.}
\end{ex}

\begin{ex}%[Thi thử TN - Chuyên Hạ Long-Quảng Ninh-L2-23]%[Duong Xuan Loi-EX6]%[2H3G3-8]
	Trong không gian cho hệ trục $Oxyz$; lấy các điểm $A(a;0;0)$, $B(0;b;0)$,  $C(0;0;c)$,\\ $D\left(a+a\sqrt{b^2+c^2};b\sqrt{a^2+c^2};c\sqrt{b^2+a^2}\right)$ với $a,b,c$ dương. Biết diện tích tam giác $ABC$ bằng $\dfrac{3}{2}$ (đơn vị diện tích) và thể tích tứ diện $ABCD$ đạt giá trị lớn nhất. Khi đó phương trình mặt phẳng $(ABD)$ là $mx+ny+pz+1=0$. Tính $m+n+p$. 
	\choice
	{$-2$}
	{\True $0$}
	{$2$}
	{$-1$}
	\loigiai{
		Phương trình mặt phẳng $(ABC)\colon\dfrac{x}{a}+\dfrac{y}{b}+\dfrac{z}{c}-1=0$.\\
		Diện tích tam giác $ABC$ bằng $S=\dfrac{\sqrt{(ab)^2+(bc)^2+(ca)^2}}{2}=\dfrac{3}{2}\Rightarrow(ab)^2+(bc)^2+(ca)^2=9$.\\
		Thể tích khối tứ diện là $V=\dfrac{1}{3}S_{ABC}\cdot h=\dfrac{h}{2}$.\\
		Với $h=\mathrm{d}\left(D,(ABC)\right)=\dfrac{\sqrt{b^2+c^2}+\sqrt{c^2+a^2}+\sqrt{a^2+b^2}}{\sqrt{\dfrac{1}{a^2}+\dfrac{1}{b^2}+\dfrac{1}{c^2}}}=abc\dfrac{\left(\sqrt{b^2+c^2}+\sqrt{c^2+a^2}+\sqrt{a^2+b^2}\right)}{3}$.\\
		$V$ lớn nhất khi $h$ lớn nhất.\\
		Ta có $abc\sqrt{b^2+c^2}=\sqrt{b^2c^2\left(a^2b^2+a^2c^2\right)}\leq\sqrt{2}\sqrt{b^2c^2\dfrac{\left(a^2b^2+a^2c^2\right)}{2}}\leq\dfrac{\sqrt{2}}{2}\left[b^2c^2+\dfrac{\left(a^2b^2+a^2c^2\right)}{2}\right]$.\\
		Tương tự $abc\sqrt{c^2+a^2}\leq\dfrac{\sqrt{2}}{2}\left[a^2c^2+\dfrac{\left(a^2b^2+b^2c^2\right)}{2}\right],abc\sqrt{b^2+a^2}\leq\dfrac{\sqrt{2}}{2}\left[a^2b^2+\dfrac{\left(a^2c^2+b^2c^2\right)}{2}\right]$.\\
		$h\leq\dfrac{2\sqrt{2}}{2\cdot 3}\left((ab)^2+(bc)^2+(ca)^2\right)=3\sqrt{2}$.\\
		Đẳng thức xảy ra khi $a=b=c=\sqrt[4]{3}$.\\
		Ta có $D\left(\sqrt[4]{3}+\sqrt{6};\sqrt{6};\sqrt{6}\right)\Rightarrow\overrightarrow{AD}=\sqrt{6}(1;1;1),\overrightarrow{AB}=\sqrt[4]{3}(-1;1;0)$, $\left[\overrightarrow{AB},\overrightarrow{AD}\right]=k(1;1;-2)$.\\
		Phương trình mặt phẳng $(ABD)\colon x+y-2z-\sqrt[4]{3}=0\Rightarrow-\dfrac{x}{\sqrt[4]{3}}-\dfrac{y}{\sqrt[4]{3}}+\dfrac{2z}{\sqrt[4]{3}}+1=0\Rightarrow\heva{&m=-\dfrac{1}{\sqrt[4]{3}}\\&n=-\dfrac{1}{\sqrt[4]{3}}\\&p=\dfrac{2}{\sqrt[4]{3}}.}$ \\
		Do đó $m+n+p=0$.}
\end{ex}
\Closesolutionfile{ans}
\begin{indapan}{10}
	{ans/ans-2-TT-19-ChuyenHaLong-QuangNinh-L2-23}
\end{indapan}